% -*- coding: utf-8 -*-
\chapter{绪\quad 论}
\label{chap:intro}

\section{课题研究背景与现实需求}

随着全球信息社会的飞速演进与万物智联时代的到来，第五代移动通信系统（5G）已全面实现商用部署。然而，面对全息通信、数字孪生、超能交通以及天地一体化泛在覆盖等面向 2030 年及未来的前沿应用场景，无线网络在峰值传输速率、空中接口时延、连接密度以及全网能效比方面提出了跨越式发展需求~\cite{wang2023road}。第六代移动通信系统（6G）应运而生，其核心愿景在于构建一个以人为中心、万物智联、通感算一体的普适智慧信息网络体系~\cite{liu2022integrated}。

为突破香农信道容量瓶颈并获取更宽广的频谱带宽资源，无线通信载波频段由微波向毫米波乃至太赫兹高频频段纵深拓展~\cite{hong2021role}。高频通信拥有数十吉赫兹的超大连续带宽，能够提供极高的数据传输吞吐量。然而，根据自由空间电磁波传播的 Friis 传输公式
\begin{equation}
  P_r = P_t G_t G_r \left( \frac{\lambda}{4\pi d} \right)^{2},
  \label{eq:friis_pathloss}
\end{equation}
高频信号波长 $\lambda$ 急剧缩短，导致自由空间路径传播损耗与载波频率的平方呈正比急剧上升。更严峻的是，高频电磁波穿透损耗极大、绕射与散射能力极度匮乏，极易受到建筑物、移动车辆甚至人体行进的物理遮挡，导致视距传播链路频繁中断，在密集城区与室内热点区域形成严重的覆盖空洞。

传统应对覆盖空洞的工程手段通常依赖于部署超密集基站或有源微型中继器。然而，有源基站与中继设备硬件结构复杂，包含昂贵的数模转换器、混频器、功率放大器等高功耗射频链路组件。超大规模高密度部署不仅带来高昂的网络建设支出与运维电费成本，更在小区间引发严重的同频同道干扰，难以满足 6G 绿色低碳可持续发展的生态底线。

在此背景下，可重构智能表面（Reconfigurable Intelligent Surface, RIS），亦称智能反射表面（Intelligent Reflecting Surface, IRS），作为一种颠覆性的新型无线物理层技术引起了学术界与产业界的广泛关注~\cite{wu2020towards}。RIS 由大量低成本、低功耗的亚波长周期性人造电磁超材料单元二维排布而成。每个反射单元通过集成 PIN 二极管或变容管，在低功耗控制器的离散控制下实时调节表面等效阻抗，从而动态改变入射电磁波的反射相位、幅度甚至偏振状态。通过智能配置超表面海量无源微元的反射相位，RIS 能够在基站与受阻用户之间协同构造虚拟可视直达链路，实现非视距向视距的电磁转折，突破了传统无线通信仅能被动适应随机衰落信道的物理局限，开辟了主动重塑电磁传播环境的全新设计范式~\cite{wu2019intelligent}。

从网络演进的视角看，RIS 的价值不仅在于覆盖补盲。已有研究指出，RIS 的定向反射能力可以显著扩展小区边缘的覆盖半径，而反射面的朝向与部署位置直接决定了可获得的阵列增益~\cite{zeng2021ris}；在硬件实现层面，超表面的电磁响应设计、偏置网络与功耗预算是决定其能否规模化商用的关键约束~\cite{rana2023ris}；在系统性能层面，非理想硬件带来的相位噪声与幅度误差会使 RIS 的理论增益出现明显折损，其折损程度与阵元规模、量化精度均强相关~\cite{boulogeorgos2020ris}。这些工作共同表明，RIS 的工程价值必须放在「硬件可实现性—信道获取代价—波束设计复杂度」三者的联合约束下评估，而不能仅停留在理想连续相移与完美信道假设下的容量分析~\cite{letaief2022edge}。

\section{可重构智能表面的物理机理与实现形态}

RIS 的核心物理基础是可编程超材料。通过在亚波长尺度上排布金属谐振单元并加载可调器件，超表面可以在宏观上呈现出等效的表面阻抗分布，进而对入射波的相位与幅度实施逐单元的精细调控。按照调控维度的不同，现有实现可分为仅调相位的相位型表面、同时调幅度与相位的幅相联合型表面，以及能够同时实现透射与反射的可重构智能表面等形态~\cite{rana2023ris}。在射频实现上，PIN 二极管提供离散的开关式相位状态，变容管提供准连续的电容调谐，二者在插入损耗、切换速度与控制复杂度上各有取舍。近年来还出现了基于液晶、微机电系统以及等离子体激元的新型调控机制，其中等离子体型超表面能够在波束偏转的同时完成极化转换，为多功能一体化提供了新的物理入口~\cite{magarotto2023ris}。

除无源反射形态外，主动式与能量选择性形态亦受到关注。主动 RIS 在反射单元中引入有源放大链路，以额外的直流功耗换取对路径损耗的补偿，从而缓解无源 RIS 在双跳级联信道中面临的双重衰落问题~\cite{ren2023ris}；能量选择性 RIS 则依据入射功率自适应地切换工作模式，在保证安全通信的同时兼顾能量效率~\cite{hong2026ris}。面向空中平台，将 RIS 搭载于无人机或高空平台可以借助三维机动性进一步优化反射几何，形成所谓的空中 RIS 与星载 RIS 形态，其可达速率与部署高度、反射单元数目之间存在复杂的耦合关系~\cite{zhou2021ris}。

\section{国内外研究现状与核心挑战}

\subsection{主被动协同波束赋形}

在 RIS 辅助的多天线基站下行多用户系统中，系统总和速率由基站端有源数字预编码与 RIS 端无源反射相移共同决定。由于反射单元的相移通常受单位模等式约束，可行域形成高度非凸的复环形流形结构，加之多用户发射信号与级联反射矩阵在信噪比表达式中深度耦合，直接联合求解属于强非凸优化难题~\cite{wu2019intelligent}。

围绕这一难题，现有研究大致沿四条技术路线展开。第一条是半正定松弛及其变体，通过矩阵升维将非凸约束松弛为半定锥约束，再借助高斯随机化恢复可行解；该方法理论完备，但当阵元数达到数百量级时，升维带来的计算复杂度会急剧膨胀，难以满足毫秒级相干时间内的实时求解要求。第二条是交替优化，将原问题分解为基站预编码子问题与 RIS 相移子问题交替求解，前者在固定相移下通常具有闭式或凸形式，后者在单位模约束下可采用流形梯度或闭式相位对齐求解。已有工作将交替优化框架推广至 MISO-OFDM 系统，联合优化子载波间的功率分配与反射波束~\cite{feng2021ris}；也有研究将贝叶斯优化引入反射波束设计，以少量函数评估逼近最优相移配置~\cite{wang2022ris}。第三条是借助额外的空间自由度或部署自由度换取性能提升，例如联合优化反射面的位置与被动波束，使等效信道条件数得到改善~\cite{guo2023ris}，或在混合波束赋形架构下同时设计射频链路与反射相位~\cite{saqib2024ris}。第四条是引入学习类方法，用神经网络或强化学习近似求解非凸映射，以离线训练代价换取在线推理的低时延~\cite{yang2021deep,yuan2024ris}。

在性能评估层面，已有研究从解析角度刻画了反射波束效率与阵元规模、入射角之间的关系，指出当反射面偏离镜面方向时阵列增益会显著下降~\cite{stratidakis2021ris}；也有工作将 RIS 用于协同干扰与安全传输，通过被动波束抑制窃听链路的信干噪比~\cite{li2022ris}，或利用空中 RIS 的机动性同时优化干扰抑制与部署位置~\cite{tang2021ris}。此外，功率分裂式有源 RIS、级联与并行混合架构下的可靠协同波束赋形等形态也被相继提出~\cite{ma2023ris,elsherbini2023ris}，表明 RIS 的波束设计已从单一相移优化扩展到「相位—幅度—功率—位置」的多维联合设计。

\subsection{级联信道状态信息获取}

主被动波束赋形的增益建立在精确信道状态信息之上。然而，出于低成本与零射频功耗的考量，无源 RIS 通常不配备基带处理芯片与射频接收机，无法直接接收或解码导频以独立估计基站—表面信道或表面—用户信道，收发两端仅能获取二者相乘构成的复合级联信道。在由 $M$ 根基站天线、$N$ 个反射单元与 $K$ 个用户构成的系统中，级联信道包含 $KMN$ 个复未知参数；采用逐个单元开关的穷举导频方案所需的最小导频符号数与 $N$ 同阶，在毫秒级时变相干块中会耗尽绝大部分传输时隙~\cite{mao2021ris}。

为压缩导频开销，现有研究主要采用三类思路。其一是结构化导频与变换域估计，利用离散傅里叶变换码本或哈达玛矩阵构造正交观测，将级联信道投影到稀疏域后以少量观测恢复~\cite{hongyun2024ris}；其二是双时间尺度估计，利用基站—表面信道变化缓慢而表面—用户信道变化较快的先验，将高维级联信道的估计分解为长周期与短周期两部分，显著降低平均导频开销~\cite{hu2021ris}；其三是借助可扩展的两阶段协议或全息式先验，把反射相移配置与信道估计联合设计，使观测矩阵随估计精度自适应调整~\cite{an2022ris}。深度学习的引入进一步降低了高维估计的在线复杂度，超分辨网络与去噪网络被用于从低分辨率观测中恢复级联信道~\cite{shen2023ris}，两阶段稀疏精化方法则通过粗估与细化的级联结构提升低信噪比下的鲁棒性~\cite{ma2026ris}。

在宽带与高速移动场景下，级联信道还表现出明显的波束斜视与多普勒扩展，已有工作将信道估计与跟踪联合建模，以卡尔曼类滤波器跟踪时变角度与增益~\cite{zhang2025ris}；全双工与硬件受限条件下的估计问题也被系统研究，指出残余自干扰与相位噪声会构成额外的估计误差底~\cite{fernandes2023ris}。此外，面向海量连接的联合活跃性检测与信道估计、以及面向太赫兹的稀疏估计方案，进一步拓展了该问题的适用边界~\cite{xia2021ris,yang2025ris}。

\subsection{离散相移量化与硬件非理想性}

实际 RIS 的相移由有限位数量的数字控制码字驱动，理想连续相移假设在实际系统中并不成立。已有研究系统分析了量化位数对可达速率的影响，指出当量化位数超过三至四位后，进一步增加位数带来的速率增益迅速饱和，而硬件成本与控制开销却线性增长~\cite{boulogeorgos2020ris}。在硬件损伤方面，相位噪声、幅度误差与单元间互耦会使波束主瓣展宽、旁瓣抬升，已有工作将硬件损伤显式建模进优化问题，通过鲁棒设计换取性能的稳定~\cite{chen2025ris}；在无小区大规模 MIMO 场景下，硬件损伤对 RIS 增益的折损程度还与导频污染强度相关~\cite{zhang2023ris}。当瞬时信道难以获取时，基于统计信道状态信息的传输设计成为一种低开销替代方案，其性能损失在用户数较多时可以被有效摊薄~\cite{dai2022ris}。

\subsection{近场、太赫兹与波束训练}

随着阵列孔径与载频的提升，RIS 辅助系统逐渐进入近场与太赫兹 regime。已有研究揭示了太赫兹频段下远场与近场波束斜视的不同成因，并给出了基于真时延与相移混合的补偿方案~\cite{hao2023ris}；在近场太赫兹通信中，极化域与距离域的自由度可被联合利用以提升波束聚焦精度~\cite{qiu2026ris}。太赫兹多用户大规模 MIMO 场景下的波束训练与混合预编码也被系统研究，指出训练开销随阵元数增长是制约其落地的核心瓶颈~\cite{ning2021ris}。为降低训练开销，分布式波束训练、基于反馈的波束训练以及联合波束训练与定位的方案相继被提出~\cite{mei2021ris,cao2023ris,wang2021ris}；面向宽带的精确波束训练与任意波束方向图快速设计方法，则从波形与码本层面进一步压缩了训练时隙~\cite{chen2023ris,dang2023ris}。深度学习类方法被用于宽带太赫兹的混合预编码设计~\cite{yuan2024ris}，而覆盖与电磁场暴露合规性的联合分析则把健康与监管约束纳入了 RIS 部署优化~\cite{melhem2025ris}。在多跳路由与多径分流场景中，主被动波束的分裂与合并策略也被用于提升端到端吞吐~\cite{mei2022ris}。

\subsection{与其他关键技术的融合}

RIS 与多址技术的结合是最活跃的分支之一。非正交多址（NOMA）方面，已有工作研究了 RIS 辅助下的用户配对与功率分配，比较了 NOMA 与正交多址在 RIS 增益下的相对优劣~\cite{zheng2020ris}，并分析了信道衰落与硬件不完美条件下的中断性能~\cite{guo2021ris,hemanth2020ris}；全双工 RIS 场景下的分层译码机制被用于抑制残余干扰~\cite{andiappan2023ris}，而主动 RIS 与无线携能 NOMA 的结合则同时考虑了速率与能量收集的折中~\cite{cantos2024ris}。速率分割多址（RSMA）方面，已有研究将 RSMA 与 RIS 联合设计以提升多用户和速率与公平性~\cite{bansal2021ris,zhang2024ris}，并将其扩展到安全传输、携能通信以及超对角 RIS 架构~\cite{gao2023ris,camana2022ris,li2024ris}。在安全与干扰管理方面，RIS 被用于构造友好干扰、抑制窃听与增强物理层安全~\cite{jiang2020ris,xiu2020ris}，非互易信道下的防御型 RIS 与能量选择性 RIS 则从硬件机理层面拓展了安全的实现路径~\cite{chenhu2025ris,hong2026ris}。在感知与通感一体化方面，RIS 辅助的双功能雷达通信系统与毫米波通感系统中的反馈式波束训练已被验证可同时改善探测与通信性能~\cite{li2023ris,cao2023ris}；基于 RIS 的混合近场远场目标定位方法则把定位精度与阵列几何的关系显式刻画出来~\cite{zhu2026ris}。在携能与物联网方向，RIS 辅助的同步无线信息与功率传输、可同时透射与反射的 STAR-RIS 无人机网络、以及低功耗广域网的物理层安全增强均有系统研究~\cite{xie2025ris,khoshafa2024ris}。

空中 RIS 是另一条重要分支。将 RIS 搭载于无人机可借助三维机动性重构反射几何，已有工作联合优化无人机轨迹与反射单元数目~\cite{zhao2026ris}、联合优化无人机位置与资源分配~\cite{nguyen2022ris}，并研究了海上无线携能通信、反向散射通信与安全通信等典型场景~\cite{yuan2025ris,wang2023ris,ye2024ris,sun2021ris}。几何随机的无人机 MIMO 信道模型与 RIS 辅助的统计信道建模也被相继提出，为轨迹与波束的联合设计提供了可信的信道基础~\cite{wang2024ris}。此外，RIS 在协作干扰、频谱共享认知网络、以及双跳 OFDM 传输中的应用亦有报道~\cite{peng2023ris,srivastava2025ris,pradhan2020ris}，而多智能反射面协同定位则把 RIS 从通信增强器件扩展为感知基础设施~\cite{ghermezchesh2023ris}。

\section{现有研究的不足与本文定位}

综合上述文献可以看出，RIS 辅助通信的研究已相当丰富，但仍存在三方面不足。

第一，波束赋形算法的高维计算开销与实际可实现性之间存在缺口。半正定松弛类方法在阵元数达到数百量级时复杂度难以承受，而多数交替优化工作仅报告理想收敛曲线，未给出在固定计算预算下的实际可达速率，也较少讨论算法在低功率区间的行为。

第二，级联信道估计的实验评估常建立在过强的先验之上。压缩感知类方法默认级联信道在变换域稀疏，但级联信道的 rank-1 逐元素调制结构会使这一先验并不总是成立；同时，部分工作在全开销条件下比较估计精度，掩盖了部分观测下的真实折中关系。

第三，离散相移量化的评估口径不统一。若量化方案与连续方案采用不同的初始点或不同的迭代预算，就可能出现量化解优于连续解的悖论，从而使量化损失被系统性低估或高估。

针对上述不足，本文以「可复现、可核验、口径一致」为原则，围绕 RIS 辅助多用户 MISO 系统的联合波束赋形与级联信道估计展开研究，主要工作包括：建立含离散相移与硬件损伤的级联信道模型；提出基于交替优化与二次变换分式规划的波束解耦机制，并给出收敛性与复杂度的形式化分析；设计基于正交 DFT 码本的部分观测级联信道估计协议，并给出与克拉美—罗下界的对照验证；在统一的计算与功率预算下系统评估量化位数、导频开销、反射单元数目与发射功率对性能的影响，并对低功率区间的用户服务数下降、压缩感知先验失配等负结果作出诚实报告。

\section{论文主要研究内容与组织结构}

本论文其余部分的组织结构安排如下：

\begin{itemize}
  \item \textbf{第二章（国内外研究现状与理论基础）}：系统回顾 RIS 电磁调控物理机理与信道传播模型，综述主被动波束赋形优化算法、级联信道估计协议与离散相移量化研究，并通过代表性文献的多维对比表分析现有工作在高维计算开销与导频冗余方面的局限。
  \item \textbf{第三章（RIS 主被动波束赋形与低开销信道估计机制）}：建立基站—表面—用户多用户 MISO 级联系统模型；给出基于交替优化与二次变换分式规划的波束解耦机制及其收敛性与复杂度分析；设计基于正交 DFT 码本的部分观测级联信道估计协议，并推导对应的克拉美—罗下界。
  \item \textbf{第四章（仿真结果与性能评估）}：给出完整的仿真参数集与评价指标，从算法收敛性、发射功率、反射单元数目、相移量化位数、导频开销与运行时延六个维度评估所提方案，并与随机相移、无 RIS、等功率 MRT 等基线在统一口径下对比，同时对负结果与适用边界作出说明。
  \item \textbf{第五章（总结与展望）}：归纳全文的研究成果，并对近场宽带效应、通感一体化、学习类方法与硬件实测等未来演进方向进行展望。
\end{itemize}
